When studying irreducibility criteria for polynomials, one question that arises is whether we can find a root in the field we are working in. For example, the polynomial $x^2-2$ has no roots in $\Q$. Could we find a way to enlarge $\Q$ so that $x^2-2$ has a root while also preserving the field structure? This leads us to the concept of a field extension. In this section, we introduce field extensions and some of their properties.

We begin by first by recalling that a field $F$ is a commutative ring with identity such that every nonzero element has an inverse. Equivalently, the set $F^\times = F \setminus \set{0}$ is an abelian group under multiplication. Associated with any field $F$ is its characteristic.
\begin{definition}[Characteristic of a Field] \leavevmode \\
    Recall that the characteristic of a field $F$, denoted $ch(F)$, is the additive order of $1$ in $F$. If the order of $1$ is infinite, we say $ch(F) = 0$.
\end{definition}

\begin{proposition}
    $ch(F) = 0$ or $p$ where $p$ is prime. If $ch(F) = p$, then
    $$\forall \alpha \in F~~p\alpha = 0$$
\end{proposition}

\begin{proof}
    If $ch(F) = 0$, we are done. Suppose $ch(F) \neq 0$. Suppose $ch(F) = n$ where $n$ is composite. Then $n = a\cdot b$ where $ 1< a < n$ and $1 < b <n$. Then 
    \begin{align*}
        (a \cdot 1_F)(b\cdot 1_F) &= \underbrace{(1_F + ... + 1_F)}_{a \text{-times}}\underbrace{(1_F + ... + 1_F)}_{b \text{-times}} \\
        &= \underbrace{1_F+...+1_F+1_F+...+1_F}_{ab \text{-times}} \\ 
        &= ab \cdot 1_F \\
        &= 0
    \end{align*}
    Since fields are integral domains, this gives $a \cdot 1_F = 0$ or $b \cdot 1_F = 0$. Either of which contradicts the minimality of $ch(F)$.
    \qed
\end{proof}

\begin{example}
    \begin{enumerate}
        \item $ch(\Q) = 0 = ch(\R)$.
        \item Denote $\Z/p\Z$ by $\F_p$. Then $ch(\F_p) = p$.
    \end{enumerate}
    In the homework, we explored the homomorphism
    \begin{align*}
        \phi&:\Z \to F \\
        \phi(n)&= n \cdot 1_F \\
    \end{align*}
    defined by 
    \begin{align*}
        &n\cdot 1 = 0 \text{ when } n = 0 \\
        &n\cdot 1 = \underbrace{1+...+1}_{n \text{-times}} \text{ when } n > 0 \\
        &n \cdot 1 = \underbrace{-1-...-1}_{n \text{-times}} \text{ when } n < 0
    \end{align*}
    $\ker \phi = n \Z$ where $n = ch(F)$. The first isomorphism theorem says
    $$\frac{\Z}{\ker \phi} \cong \phi(\Z)$$
    That is,
    $$\frac{\Z}{ch(F)} \cong \phi(\Z) \subseteq F$$
    Then when $ch(F) = 0$ we have
    $$Z \cong \frac{\Z}{\set{0}} \cong \phi(\Z)$$
    so a field of characteristic $0$ contains an isomorphic copy of $\Z$ and thus also contains an isomorphic copy of $\Q$. In the case that $ch(F) = p$ a prime, we have
    $$\frac{\Z}{p\Z}\cong \phi(\Z)$$
    Hence any field with characteristic $p$ a prime contains an isomorphic copy of $\Z/p\Z$.
\end{example}

\begin{definition}[Prime Subfield] \leavevmode \\
    The prime subfield of a field $F$ is the subfield of $F$ generated by the multiplicative identity of $F$, $1_F$. It is isomorphic to $\Q$ (if $ch(F) = 0$) or $\F_p$ (if $ch(F) = p$ a prime).
\end{definition}

\begin{definition}[Extension Field] \leavevmode \\
    If $K$ is a field containing the subfield $F$, then $K$ is said to be an extension field of $F$ (or an extension of $F$) denoted $K/F$ (read $K$ over $F$, NOT $K \mod F$), or by the diagram:
    $$
    \begin{tikzcd}[row sep=0.5cm]
        K \arrow [d,-] \\
        F
    \end{tikzcd}
    $$
    In particular, every field is an extension of its prime subfield.
\end{definition}

If $K/F$ is any extension of fields, then $K$ is a vector space over $F$. Note that $\langle K, + \rangle$ is an abelian group and since $F\subseteq K$ the scalar multiplication is defined by the multiplication in $K$. In particular, every field $F$ can be considered a vector space over its prime subfield.

\begin{definition}[Degree of a Field Extension] \leavevmode \\
    The degree (or relative degree or index) of a field extension $K/F$, denoted $[K:F]$, is the dimension of $K$ as a vector space over $F$. The extension is said to be finite if $[K:F]$ is finite, otherwise we say the extension is infinite.
\end{definition}

Our first big result about field extensions: If we have an irreducible polynomial $p(x) \in F[x]$ we can find a field $K$ containing $F$ where $p(x)$ has a zero.

\begin{recall}
    For $\phi:F \to F'$ a homomorphism of fields we have $\phi$ is either identically zero or $\phi$ is injective. So
    $$\phi(F) = 0 \text{ or } \phi(F) \cong F$$
\end{recall}

\begin{proposition}
    Let $F$ be a field and let $p(x) \in F[x]$ be an irreducible polynomial. Then there exists a field $K$ containing an isomorphic copy of $F$ in which $p(x)$ has a root.
\end{proposition}

\begin{proof}
    Consider the quotient
    $$K = \frac{F[x]}{\left(p(x)\right)}$$
    We have discovered that $\left(p(x)\right)$ is maximal and $K$ is a field. Letting $\phi = \pi \vert_F : F \to K$ (where $\pi$ is the natural projection from $F[x]$ to $F[x]/\left(p(x)\right)$ restricted to the constant polynomials in $F[x]$, where $\phi(f) = \pi(f) = f + \left(p(x)\right) ~~\forall f \in F$) gives a homomorphism of fields. Note that this mapping is not identically zero since $\phi(1_F) = 1_F + \left(p(x)\right) \neq 0 + \left(p(x)\right)$ since $K$ is a field and $1_F \neq 0$. Hence
    $$\phi(F) \cong F$$
    Therefore $\phi(F)$ is our isomorphic copy of $F$ in $K$. We now consider $F$ a subfield of $K$. Note then that there is an isomorphic copy of $F[x]$ contained in $K[x]$. Let $\bar{x} = \pi(x) = x+\left(p(x)\right) \in K$. Then $p(\bar{x}) = \bar{p(x)}$ since $\pi$ is a homomorphism. Then
    $$p(\bar{x}) = p(x) + \left(p(x)\right) = 0 + \left(p(x)\right)$$
    so $\bar{x} \in K$ and $\bar{x}$ is the root of $p(x)$. Hence, $K$ is an extension of $F$ in which $p(x)$ has a root.
    \qed
\end{proof}

\begin{theorem}
    Let $p(x) \in F[x]$ be an irreducible polynomial of degree $n$. Let $\theta = x \mod p(x)$ in $K$. Then the elements
    $$1, \theta, \theta^2, ..., \theta^{n-1}$$
    form a basis for $K$ as a vector space over $F$ so the degree of the extension $K/F$ is $n$, i.e. $[K:F] = n$. So $K = \set{a_0 + a_1\theta + a_2\theta^2 +... + a_{n-1}\theta^{n-1} : a_0, ..., a_{n-1}\in F}$ consists of all polynomials of degree less than $n$.
\end{theorem}